\documentclass{amsart}

% 数学相关宏包
\usepackage{amsmath,mathtools,amsthm,amssymb}
\usepackage{bbm}
\usepackage{mathrsfs}
\usepackage{centernot}

% 图形和绘图
\usepackage{graphicx,float}
\usepackage{tikz,tikz-cd}
\usepackage{pgfplots}
\usepackage[framemethod=TikZ]{mdframed}
\usetikzlibrary{decorations.pathmorphing,positioning,arrows,automata,calc,shapes,patterns}
\pgfplotsset{compat=1.18}
\usepgfplotslibrary{statistics}

% 表格和列表
\usepackage{booktabs,multirow,colortbl}
\usepackage{enumitem}
\setlist[itemize,enumerate,description]{noitemsep}
\setlistdepth{9}

% 算法
\usepackage{algorithm,algorithmic}

% 字体和颜色
\usepackage{xcolor,listings}
\usepackage{xeCJK}
\setCJKmainfont{Noto Serif CJK SC}

% 页面和链接
\usepackage{fancyhdr,caption}
\usepackage{hyperref}
\usepackage{cleveref}

% 工具宏包
\usepackage{etoolbox,letltxmacro,docmute}
\usepackage{gensymb}

% 自定义设置
\renewcommand{\arraystretch}{1.5}
\let\originalmathbb\mathbb
\renewcommand{\mathbb}[1]{%
    \ifstrequal{#1}{1}{\mathbbm{#1}}{\originalmathbb{#1}}%
}

% 颜色定义
\definecolor{lightblue}{RGB}{173,216,230}
\definecolor{lightgreen}{RGB}{144,238,144}
\definecolor{lightyellow}{RGB}{255,255,224}
\definecolor{lightgray}{RGB}{211,211,211}

% 超链接设置
\hypersetup{
    colorlinks=true,
    linkcolor=blue,
    filecolor=magenta,
    urlcolor=cyan,
    pdftitle={\@title},
    pdfauthor={\@author},
    pdfsubject={数学笔记},
    pdfkeywords={数学},
    bookmarksnumbered=true,
    bookmarksopen=true,
    bookmarksopenlevel=6,
    pdfstartview=FitH,
    pdfpagemode=UseOutlines,
    pdfborder={0 0 0}
}

% 定理环境
\theoremstyle{plain}
\newtheorem{theorem}{Theorem}[section]
\newtheorem{lemma}[theorem]{Lemma}
\newtheorem{corollary}[theorem]{Corollary}
\newtheorem{proposition}[theorem]{Proposition}

\theoremstyle{definition}
\newtheorem{definition}[theorem]{Definition}
\newtheorem{example}[theorem]{Example}
\newtheorem{exercise}[theorem]{Exercise}
\newtheorem{problem}[theorem]{Problem}

\theoremstyle{remark}
\newtheorem{remark}[theorem]{Remark}
\newtheorem{note}[theorem]{Note}
\newtheorem{observation}[theorem]{Observation}

% 解答环境
\newtheorem*{solution}{Solution}
\newtheorem*{proof*}{Proof}

% 图片路径
\graphicspath{{F:/iCloudDrive/iCloud~md~obsidian/2025-summer/figures/}}

\author{乐绎华, Yue Yihua}
\address{中山大学, Sun Yat-sen University}
\email{yueyh@mail2.sysu.edu.cn}
\date{\today}
\title{Mathematical Notes}

\begin{document}

\maketitle

\begin{figure}[H]
\centering
\includegraphics[width=\textwidth]{inner-product-2025070721.png}
% \caption{}
\label{}
\end{figure}

We need to show that

\begin{itemize}
	\item The conjugate symmetry: $\left< p,q \right>=\overline{\left< q,p \right>}$.
	\item Linearity in the first argument: $\left< (a\odot p)\oplus (b\odot q),r \right>=a\left< p,r \right>+b\left< q,r \right>$.
	\item Prositive definiteness: if $p$ is nonzero, then $\left< p,p \right>>0$.
\end{itemize}

Clearly,
\[
\left< p,q \right> =\sum_{i=0}^{n} \text{clr}_i(p)\text{clr}_i(q)=\sum_{i=0}^{n} \text{clr}_i(q)\text{clr}_i(p)= \underbrace{ \left< q,p \right> }_{ \text{is real} } =\overline{\left< q,p \right> }
\]
And
\[
\begin{aligned}
\sum_{i=0}^{n} \text{clr}_i(p)\text{clr}_i(q) & =\sum_{i=0}^{n} \log\frac{p_i}{\left( \prod_{j=0}^{n} p_j \right)^{1/(n+1)}}\cdot \log\frac{q_i}{\left( \prod_{j=0}^{n} q_j \right)^{1/(n+1)}} \\
 & =\sum_{i=0}^{n} \left( \log p_i-\frac{1}{n+1}\sum_{j=0}^{n} \log p_j \right)\left( \log q_i-\frac{1}{n+1}\sum_{j=0}^{n} \log q_j \right) \\
 & =\sum_{i=0}^{n} \log p_i\log q_i -\frac{1}{n+1}\left( \sum_{i=0}^{n} \log p_i \right)\left( \sum_{i=0}^{n} \log q_i \right)
\end{aligned}
\]
\[
\begin{aligned}
\frac{1}{2(n+1)}\sum_{i=0}^{n} \sum_{j=0}^{n} \log\frac{p_i}{p_j}\log\frac{q_i}{q_j} & =\frac{1}{2(n+1)}\sum_{i=0}^{n} \sum_{j=0}^{n} (\log p_i-\log p_j)(\log q_i-\log q_j) \\
 & =\frac{1}{2(n+1)}\sum_{i=0}^{n} \sum_{j=0}^{n} (\log p_i\log q_i-\log p_i\log q_j-\log p_j\log q_i+\log p_j\log q_j) \\
 & =\sum_{i=0}^{n} \log p_i\log q_i-\frac{1}{n+1}\left( \sum_{i=0}^{n} \log p_i \right)\left( \sum_{i=0}^{n} \log q_i \right) \\
 & =\sum_{i=0}^{n} \text{clr}_i(p)\text{clr}_i(q)
\end{aligned}
\]
Then
\[
\begin{aligned}
(a\odot p)\oplus (b\odot q) & =\left[ \frac{(p_0)^{a}}{\sum_{i=0}^{n} (p_i)^{a}},\dots,\frac{(p_n)^{a}}{\sum_{i=0}^{n} (p_i)^{a}} \right]\oplus \left[ \frac{(q_0)^{b}}{\sum_{i=0}^{n} (q_i)^{b}},\dots,\frac{(q_n)^{b}}{\sum_{i=0}^{n} (q_i)^{b}} \right] \\
 & =\left[ \frac{\frac{(p_0)^{a}}{\sum_{i=0}^{n} (p_i)^{a}}\cdot\frac{(q_0)^{b}}{\sum_{i=0}^{n} (q_i)^{b}}}{\sum_{j=0}^{n} \left[ \frac{(p_j)^{a}}{\sum_{i=0}^{n} (p_i)^{a}}\cdot\frac{(q_j)^{a}}{\sum_{i=0}^{n} (q_i)^{a}} \right]},\dots, \frac{\frac{(p_n)^{a}}{\sum_{i=0}^{n} (p_i)^{a}}\cdot\frac{(q_n)^{b}}{\sum_{i=0}^{n} (q_i)^{b}}}{\sum_{j=0}^{n} \left[ \frac{(p_j)^{a}}{\sum_{i=0}^{n} (p_i)^{a}}\cdot\frac{(q_j)^{a}}{\sum_{i=0}^{n} (q_i)^{a}} \right]}\right]
\end{aligned}
\]
Thus
\[
\begin{aligned}
\left< (a\odot p)\oplus (b\odot q),r \right> & =\frac{1}{2(n+1)}\sum_{i=0}^{n} \sum_{j=0}^{n} \log\frac{(p_i)^{a}\cdot(q_i)^{b}}{(p_j)^{a}\cdot(q_j)^{b}}\log\frac{r_i}{r_j} \\
 & =\frac{1}{2(n+1)}\sum_{i=0}^{n} \sum_{j=0}^{n} \left( a\log\frac{p_i}{p_j}+b\log\frac{q_i}{q_j} \right)\log\frac{r_i}{r_j} \\
 & =a\left< p,r \right> +b\left< q,r \right> 
\end{aligned}
\]
And for nonzero $p\in\Delta _n$ (nonzero means $p\neq o$, where $o=\left( \frac{1}{n+1},\dots,\frac{1}{n+1} \right)$)
\[
\left< p,p \right> =\frac{1}{2(n+1)}\sum_{i=0}^{n} \sum_{j=0}^{n} \left( \log\frac{p_i}{p_j} \right)^{2}\geq 0
\]
The equality holds only when $p_i=p_j,\forall i,j$, i.e. $p=o$. Then for $p\neq o$,
\[
\left< p,p \right> >0
\]
Hence $\left< \cdot,\cdot  \right>$ is an inner product.


\end{document}